\documentclass{templates/forecastp1}

\newcommand{\authorinfo}{THE FORUM CENTER - NGUYỄN HOÀNG HUY}
\newcommand{\documenttitle}{SPEAKING FORECAST QUÝ III}
\newcommand{\collectiontitle}{TUYỂN TẬP SAMPLE SPEAKING IELTS FORECAST QUÝ III (P1)}

\begin{document}

\thispagestyle{firstpage}
\makeforecastheader{\authorinfo}{\documenttitle}{\collectiontitle}
\pagestyle{otherpages}

\forecasttopic{Topic 1: Pets and Animals}
\forecastqa{What's your favourite animal? Why?}{I'd have to say dogs are my absolute favourite. There's something incredibly comforting about their loyalty and how they can instantly lift your spirits after a stressful day of lectures. I've always been drawn to their playful nature, and honestly, watching dog videos during study breaks has become my guilty pleasure.}
\forecastqa{Where do you prefer to keep your pet, indoors or outdoors?}{Definitely indoors, especially in Vietnam's climate. I think pets, particularly dogs and cats, thrive better when they're part of the household. Plus, keeping them inside means they're safer from traffic and diseases, which is quite a concern in urban areas like where I live. It also makes bonding with them much easier.}
\forecastqa{Have you ever had a pet before?}{Yes, I had a golden retriever named Lucky throughout my childhood. He was basically my best friend growing up – we'd play together every afternoon after school. Sadly, he passed away just before I started university, which was absolutely heartbreaking. I still miss him terribly, but those memories are priceless.}
\forecastqa{What is the most popular animal in your country?}{I'd say dogs are probably the most popular pets in Vietnam, though cats are catching up fast, especially among young people in cities. Traditionally, many Vietnamese families also keep fish for feng shui purposes. In rural areas, you'll still see chickens and pigs, but those are mainly for practical reasons rather than companionship.}
\textbf{Vocabulary}\par
\begin{itemize}
  \item \textbf{Lift your spirits} -- to make you feel happier and more positive -- làm bạn vui vẻ, phấn chấn hơn
  \item \textbf{Guilty pleasure} -- something you enjoy despite feeling you shouldn't -- thú vui tội lỗi, sở thích ngại ngùng thừa nhận
  \item \textbf{Thrive} -- to grow or develop successfully -- phát triển tốt, khỏe mạnh
  \item \textbf{Heartbreaking} -- extremely sad and upsetting -- vô cùng đau lòng
  \item \textbf{Catching up} -- becoming as popular or successful as something else -- đang theo kịp, bắt kịp
  \item \textbf{Feng shui purposes} -- for traditional Chinese beliefs about harmony and energy -- mục đích phong thủy
  \item \textbf{Companionship} -- friendship and company -- tình bạn, sự đồng hành
\end{itemize}

\forecasttopic{Topic 2: Dreams}
\forecastqa{Can you remember the dreams you had?}{Occasionally, yes, but they tend to fade pretty quickly once I wake up. I find that the more vivid ones usually stick with me – like last week I dreamt I was late for an exam, which is probably my subconscious reflecting my exam anxiety. The mundane dreams just slip away within minutes of waking up.}
\forecastqa{Do you share your dreams with others?}{Not really, unless they're particularly bizarre or funny. I'll sometimes tell my roommate about weird dreams over breakfast, and we'll have a good laugh about them. But generally, I keep them to myself since most dreams don't make much sense when you try to explain them out loud.}
\forecastqa{Do you think dreams have special meanings?}{I'm quite skeptical about that, to be honest. While I know some cultures place great emphasis on dream interpretation, I tend to think they're just our brain's way of processing information and emotions from the day. That said, recurring dreams might indicate underlying stress or concerns that we should pay attention to.}
\forecastqa{Do you want to make your dreams come true?}{If you're talking about aspirations rather than sleep dreams, then absolutely! I'm working towards my dream of studying abroad for a master's degree, which requires a lot of preparation and saving up. It's challenging, but having that goal keeps me motivated through tough times at university.}
\textbf{Vocabulary}\par
\begin{itemize}
  \item \textbf{Vivid} -- producing clear, powerful mental images -- sống động, rõ nét
  \item \textbf{Subconscious} -- the part of your mind you're not fully aware of -- tiềm thức
  \item \textbf{Mundane} -- ordinary and not interesting -- tầm thường, bình thường
  \item \textbf{Bizarre} -- very strange or unusual -- kỳ quặc, lạ lùng
  \item \textbf{Skeptical} -- having doubts about something -- hoài nghi, nghi ngờ
  \item \textbf{Place great emphasis on} -- to consider something very important -- đặt nặng vấn đề, coi trọng
  \item \textbf{Underlying} -- hidden but important -- tiềm ẩn, cơ bản
\end{itemize}

\forecasttopic{Topic 3: Morning time}
\forecastqa{Do you like getting up early in the morning?}{Honestly, I'm not naturally a morning person at all. I find it quite challenging to drag myself out of bed before 7 AM, especially during winter when it's still dark outside. However, I've noticed that on the rare occasions when I do wake up early, I feel more productive throughout the day.}
\forecastqa{What do you usually do in the morning?}{My morning routine is pretty straightforward. I typically wake up around 6:30, scroll through social media for a few minutes – which I know is a bad habit – then freshen up and grab a quick breakfast. If I have an 8 AM lecture, I'll just have some bread and coffee before rushing out the door.}
\forecastqa{What did you do in the morning when you were little? Why?}{Back then, mornings were much more leisurely. My mum would wake me up, and I'd have a proper sit-down breakfast with my family – usually phở or rice porridge. Then I'd watch some cartoons before heading to school. I had much more time because my primary school was just around the corner from home.}
\forecastqa{Are there any differences between what you do in the morning now and what you did in the past?}{Absolutely, there's a world of difference. These days, everything feels more rushed and hectic. I've had to become much more independent with getting ready, and breakfast is often something I grab on the go rather than a family affair. The biggest change is probably the stress level – mornings now involve checking deadlines and assignment due dates.}
\forecastqa{Do you spend your mornings doing the same things on both weekends and weekdays? Why?}{Not at all. Weekday mornings are all about efficiency and getting to class on time, whereas weekends are my chance to actually sleep in until 9 or 10. On Saturdays and Sundays, I take my time with breakfast, maybe catch up on a TV series, and just enjoy a slower pace without the pressure of lectures looming.}
\textbf{Vocabulary}\par
\begin{itemize}
  \item \textbf{Drag myself out of bed} -- to force yourself to get up when you don't want to -- ép buộc bản thân thức dậy
  \item \textbf{Straightforward} -- simple and easy to understand -- đơn giản, rõ ràng
  \item \textbf{Freshen up} -- to wash and make yourself clean and tidy -- rửa mặt, sửa soạn
  \item \textbf{Leisurely} -- relaxed and unhurried -- thong thả, không vội vã
  \item \textbf{A world of difference} -- a very big difference -- khác biệt rất lớn
  \item \textbf{Hectic} -- very busy and full of activity -- bận rộn, hối hả
  \item \textbf{On the go} -- while moving or traveling -- khi đang di chuyển
  \item \textbf{Looming} -- appearing as a threat or difficulty -- đang đến gần, lờ mờ đe dọa
\end{itemize}

\forecasttopic{Topic 4: Key}
\forecastqa{Do you always bring a lot of keys with you?}{Not really, I try to keep it minimal. I usually just carry my dorm room key and my motorbike key on a small keychain. I used to have a bunch of keys for different things, but I found it quite cumbersome, so I've streamlined it to just the essentials.}
\forecastqa{Have you ever lost your keys?}{Unfortunately, yes, and it was a nightmare. Last semester, I misplaced my dorm key after a late-night study session at the library. I had to wake up my roommate at 2 AM to let me in, which was incredibly embarrassing. Since then, I've been much more careful and always double-check my pockets.}
\forecastqa{Do you often forget the keys and lock yourself out?}{It's happened a couple of times, I'll admit. Usually when I'm rushing out for an early class and my mind is elsewhere. The most frustrating part is having to wait around for my roommate to come back, especially if I've got assignments to work on. Now I've developed a habit of patting my pocket before closing the door.}
\forecastqa{Do you think it's a good idea to leave your keys with a neighbour?}{I'm a bit torn on this one. On one hand, it's practical to have a backup option in case of emergencies. On the other hand, I'm quite cautious about security, especially in a dorm setting where you don't always know your neighbors well. I'd probably only do it with someone I really trust, like a close friend living nearby.}
\textbf{Vocabulary}\par
\begin{itemize}
  \item \textbf{Cumbersome} -- heavy, complicated, or difficult to use -- cồng kềnh, rườm rà
  \item \textbf{Streamlined} -- made simpler and more efficient -- đơn giản hóa, tinh gọn
  \item \textbf{A nightmare} -- a very difficult or unpleasant experience -- cơn ác mộng, trải nghiệm tồi tệ
  \item \textbf{Incredibly embarrassing} -- extremely causing shame or awkwardness -- vô cùng xấu hổ
  \item \textbf{Torn on} -- unable to decide between two options -- phân vân, do dự
  \item \textbf{Cautious about} -- careful to avoid problems or danger -- thận trọng về
  \item \textbf{Backup option} -- an alternative plan or solution -- phương án dự phòng
\end{itemize}

\forecasttopic{Topic 5: Day off}
\forecastqa{When was the last time you had a few days off?}{That would be during the Lunar New Year holiday, which was about two months ago. I had nearly a week off from university, which felt like absolute bliss after a grueling exam period. It was the perfect opportunity to recharge my batteries and spend quality time with my family back in my hometown.}
\forecastqa{What do you usually do when you have days off?}{It really depends on my mood and energy levels. Sometimes I'll use the time to catch up on sleep and binge-watch Netflix series I've been meaning to see. Other times, I'll meet up with friends for coffee or explore new cafes around the city. If I'm feeling productive, I might tackle some personal projects or read books for pleasure rather than for assignments.}
\forecastqa{Do you usually spend your days off with your parents or with your friends?}{It's quite split, actually. During longer breaks like Tet, I'll definitely head home to spend time with my parents and extended family – that's non-negotiable in Vietnamese culture. But for shorter breaks or weekends, I tend to hang out with my university friends since we're all in the same city and can easily coordinate plans.}
\forecastqa{What would you like to do if you had a day off tomorrow?}{Oh, that would be fantastic! I'd probably start with a proper lie-in until late morning, then treat myself to brunch at this cozy cafe I've been wanting to try. In the afternoon, I might visit a bookstore or art gallery, and in the evening, catch up with some friends over dinner. Basically, anything that doesn't involve textbooks or lecture halls!}
\textbf{Vocabulary}\par
\begin{itemize}
  \item \textbf{Absolute bliss} -- complete happiness -- hạnh phúc tuyệt đối
  \item \textbf{Grueling} -- extremely tiring and demanding -- mệt mỏi, khắc nghiệt
  \item \textbf{Recharge my batteries} -- to rest and regain energy -- nạp lại năng lượng
  \item \textbf{Binge-watch} -- to watch many episodes of a TV series in one sitting -- xem liền tù tì nhiều tập
  \item \textbf{Non-negotiable} -- not open to discussion or change -- không thể thương lượng, bắt buộc
  \item \textbf{Coordinate plans} -- to organize activities together -- sắp xếp kế hoạch cùng nhau
  \item \textbf{Lie-in} -- staying in bed later than usual -- nằm dài trên giường, ngủ nướng
\end{itemize}

\forecasttopic{Topic 6: Gifts}
\forecastqa{What gift have you received recently?}{Just last month, my best friend gave me a really thoughtful gift for my birthday – a Kindle e-reader. She knows I'm an avid reader but always complain about lugging heavy textbooks around campus. It was such a practical yet personal present, and I've been using it constantly ever since.}
\forecastqa{Have you ever sent handmade gifts to others?}{Yes, actually! Last Christmas, I made personalized photo albums for my close friends, filled with pictures from our university adventures together. It was quite time-consuming to put together, but I wanted to give them something meaningful rather than just buying generic presents. The reactions I got made all the effort worthwhile.}
\forecastqa{Have you ever received a great gift?}{Definitely – the Kindle I mentioned earlier is probably the best gift I've received in recent years. But if I'm being nostalgic, the guitar my parents gave me when I turned 16 holds a special place in my heart. It sparked my interest in music, and I still play it regularly to unwind after stressful study sessions.}
\forecastqa{What do you consider when choosing a gift?}{I always try to think about the person's interests and what would genuinely be useful or bring them joy. Budget is obviously a factor as a student, but I believe thoughtfulness matters more than price tags. I also consider whether it's something they wouldn't typically buy for themselves – that element of surprise makes it more special.}
\forecastqa{Do you think you are good at choosing gifts?}{I'd like to think so! I tend to pay attention to little hints people drop in conversations about things they want or need. My friends usually seem genuinely pleased with what I give them, so I must be doing something right. That said, I still occasionally second-guess myself and worry whether they'll actually like it.}
\textbf{Vocabulary}\par
\begin{itemize}
  \item \textbf{Avid reader} -- someone who reads a lot and enthusiastically -- người ham đọc sách
  \item \textbf{Lugging} -- carrying something heavy with difficulty -- kéo lê, mang vác nặng nhọc
  \item \textbf{Time-consuming} -- taking a lot of time -- tốn thời gian
  \item \textbf{Worthwhile} -- worth the time, effort, or money spent -- đáng giá, xứng đáng
  \item \textbf{Nostalgic} -- feeling sentimental about the past -- hoài niệm, nhớ về quá khứ
  \item \textbf{Unwind} -- to relax after stress or tension -- thư giãn, xả stress
  \item \textbf{Second-guess myself} -- to doubt your own decisions -- nghi ngờ quyết định của bản thân
\end{itemize}

\forecasttopic{Topic 7: Hobby}
\forecastqa{Do you have any hobbies?}{Yes, I've got a few that I try to make time for despite my hectic university schedule. Photography is probably my main passion – I love capturing street scenes and candid moments around Hanoi. I also enjoy playing guitar and occasionally dabble in creative writing when inspiration strikes, though that's become less frequent with all the academic essays I have to write.}
\forecastqa{Did you have any hobbies when you were a child?}{I was really into drawing and painting back then. I'd spend hours sketching cartoon characters and creating colorful artwork, much to my parents' delight. I also collected stamps for a while, which seems quite old-fashioned now, but it taught me patience and attention to detail. Swimming was another big one – I practically lived at the pool during summer holidays.}
\forecastqa{Do you have a hobby that you've had since childhood?}{Not consistently, to be honest. While I still occasionally pick up a pencil to sketch, it's nowhere near as frequent as when I was younger. The guitar is probably the closest thing, as I started learning in high school and have stuck with it. Most of my childhood hobbies have evolved or been replaced by new interests as I've grown older.}
\forecastqa{Do you have the same hobbies as your family members?}{There's some overlap, actually. My dad and I both enjoy photography, though he's more into landscape photography while I prefer urban and street photography. My mum loves reading, which is something we bond over – we often recommend books to each other. My younger brother is into gaming, which I occasionally join him in, though I'm nowhere near his skill level.}
\textbf{Vocabulary}\par
\begin{itemize}
  \item \textbf{Hectic} -- very busy and full of activity -- bận rộn, hối hả
  \item \textbf{Candid moments} -- natural, unposed moments -- khoảnh khắc tự nhiên, không dàn dựng
  \item \textbf{Dabble in} -- to do something occasionally without serious commitment -- làm thử, nghịch ngợm với
  \item \textbf{Inspiration strikes} -- when you suddenly feel creative or motivated -- khi cảm hứng ập đến
  \item \textbf{Old-fashioned} -- not modern, belonging to the past -- lỗi thời, cổ điển
  \item \textbf{Overlap} -- to share common features or interests -- trùng lặp, có điểm chung
  \item \textbf{Bond over} -- to form a connection through shared interests -- gắn kết thông qua
\end{itemize}

\forecasttopic{Topic 8: Team sports}
\forecastqa{Do you like team sports?}{I have mixed feelings about them, honestly. I appreciate the camaraderie and social aspect of team sports – there's something energizing about working together towards a common goal. However, I'm not particularly athletic, so I sometimes feel like I'm letting my teammates down. I prefer watching team sports rather than actively participating in competitive matches.}
\forecastqa{Have you ever played a team sport?}{Yes, I played volleyball quite regularly during high school as part of our PE classes. It was compulsory, so I didn't have much choice! I was never the star player by any means, but I enjoyed the experience and the friendships that came with it. At university, I occasionally join casual football matches with my dorm mates, though nothing too serious.}
\forecastqa{What benefits do team sports bring?}{There are loads of benefits, I think. Obviously, there's the physical fitness aspect, but beyond that, team sports teach you valuable life skills like communication, cooperation, and handling pressure. They also help build lasting friendships and provide a healthy outlet for stress. For students especially, they're a great way to take a break from academic pressures and stay active.}
\forecastqa{Do you prefer individual sports or team sports?}{I'd lean towards individual sports, if I'm being honest. I enjoy the flexibility and independence they offer – you can practice at your own pace without coordinating with others. Activities like running, swimming, or cycling appeal to me more because there's less pressure and I can focus on personal improvement rather than worrying about team performance. That said, I do recognize the unique social benefits that team sports provide.}
\textbf{Vocabulary}\par
\begin{itemize}
  \item \textbf{Mixed feelings} -- having both positive and negative emotions about something -- cảm xúc lẫn lộn, phức tạp
  \item \textbf{Camaraderie} -- friendship and trust among people in a group -- tình bạn, tình đồng đội
  \item \textbf{Letting down} -- disappointing someone -- làm ai đó thất vọng
  \item \textbf{By any means} -- in any way; at all -- theo bất kỳ cách nào, chút nào
  \item \textbf{Loads of} -- a lot of -- rất nhiều
  \item \textbf{Outlet for stress} -- a way to release tension or pressure -- lối thoát cho căng thẳng
  \item \textbf{Lean towards} -- to prefer or tend to choose -- nghiêng về, thiên về
\end{itemize}

\end{document}
